\documentclass{article}
\usepackage{amsmath, amssymb, amsthm}
\usepackage[utf8]{inputenc}
\usepackage[margin=1in]{geometry}

\theoremstyle{definition}
\newtheorem{problem}{Problem}
\newtheorem*{solution}{Solution}

\title{MA 408: Measure Theory - Tutorial 1 Solutions}
\author{}
\date{}

\begin{document}

\maketitle

\begin{problem}
Suppose $X, Y$ are nonempty sets such that $X \subseteq Y \subseteq \mathcal{P}(\Omega)$ for some nonempty set $\Omega$. Show that $\mathcal{A}(X) \subseteq \mathcal{A}(Y)$ and $\mathcal{F}(X) \subseteq \mathcal{F}(Y)$. In particular, show that $\mathcal{F}(\mathcal{A}(X)) = \mathcal{F}(X)$.
\end{problem}

\begin{solution}
\textbf{1. Monotonicity of Generated Algebras and $\sigma$-Algebras}

Let $\mathcal{A}(S)$ denote the algebra generated by a collection $S$, and $\mathcal{F}(S)$ denote the $\sigma$-algebra generated by $S$. By definition:
\begin{itemize}
    \item $\mathcal{A}(S) = \bigcap \{ \mathcal{G} \mid \mathcal{G} \text{ is an algebra on } \Omega \text{ and } S \subseteq \mathcal{G} \}$.
    \item $\mathcal{F}(S) = \bigcap \{ \Sigma \mid \Sigma \text{ is a } \sigma\text{-algebra on } \Omega \text{ and } S \subseteq \Sigma \}$.
\end{itemize}

\textit{Proof of $\mathcal{A}(X) \subseteq \mathcal{A}(Y)$:}
Given $X \subseteq Y$, and observing that $Y \subseteq \mathcal{A}(Y)$ (since a generated algebra contains its generators), we have:
\[ X \subseteq Y \subseteq \mathcal{A}(Y) \implies X \subseteq \mathcal{A}(Y) \]
Since $\mathcal{A}(Y)$ is an algebra containing $X$, and $\mathcal{A}(X)$ is the smallest algebra containing $X$ (intersection of all such algebras), it implies:
\[ \mathcal{A}(X) \subseteq \mathcal{A}(Y) \]

\textit{Proof of $\mathcal{F}(X) \subseteq \mathcal{F}(Y)$:}
Similarly, since $Y \subseteq \mathcal{F}(Y)$, we have $X \subseteq \mathcal{F}(Y)$. Since $\mathcal{F}(Y)$ is a $\sigma$-algebra containing $X$, and $\mathcal{F}(X)$ is the smallest such $\sigma$-algebra:
\[ \mathcal{F}(X) \subseteq \mathcal{F}(Y) \]

\textbf{2. Equality of $\mathcal{F}(\mathcal{A}(X))$ and $\mathcal{F}(X)$}

We prove equality by showing mutual inclusion.

$(\supseteq)$ Direction:
Since $X \subseteq \mathcal{A}(X)$, we can apply the monotonicity of $\mathcal{F}$ derived above (replacing the subset $X$ with $X$ and superset $Y$ with $\mathcal{A}(X)$):
\[ \mathcal{F}(X) \subseteq \mathcal{F}(\mathcal{A}(X)) \]

$(\subseteq)$ Direction:
Note that $\mathcal{F}(X)$ is a $\sigma$-algebra containing $X$. Since every $\sigma$-algebra is also an algebra, $\mathcal{F}(X)$ is an algebra containing $X$.
Since $\mathcal{A}(X)$ is the smallest algebra containing $X$, we must have:
\[ \mathcal{A}(X) \subseteq \mathcal{F}(X) \]
Now, consider $\mathcal{F}(\mathcal{A}(X))$. This is the smallest $\sigma$-algebra containing $\mathcal{A}(X)$. Since $\mathcal{F}(X)$ is a $\sigma$-algebra that contains $\mathcal{A}(X)$ (as shown immediately above), it follows that:
\[ \mathcal{F}(\mathcal{A}(X)) \subseteq \mathcal{F}(X) \]

\noindent Since both inclusions hold, $\mathcal{F}(\mathcal{A}(X)) = \mathcal{F}(X)$.
\end{solution}

\hrulefill

\begin{problem}
Consider $\Omega = \mathbb{R}$ and consider the semi-algebra $\mathcal{C} = \{(a, b] \mid -\infty \le a \le b \le \infty\}$ on $\mathbb{R}$. Define the length function $\lambda: \mathcal{C} \to [0, \infty]$ by:
\[ \lambda((a,b]) = \begin{cases} \infty & \text{if } a = -\infty \\ b-a & \text{if } a,b \in \mathbb{R} \\ \infty & \text{if } b = \infty \end{cases} \]
Show that $\lambda$ is a measure on $\mathcal{C}$.
\end{problem}

\begin{solution}
To prove $\lambda$ is a measure on the semi-algebra $\mathcal{C}$, we must verify three properties:
\begin{enumerate}
    \item $\lambda(\emptyset) = 0$.
    \item Finite additivity.
    \item Countable additivity ($\sigma$-additivity).
\end{enumerate}

\textbf{1. Empty Set:}
The empty set corresponds to intervals where $a=b$ (for $a \in \mathbb{R}$).
\[ \lambda(\emptyset) = \lambda((a,a]) = a - a = 0 \]

\textbf{2. Finite Additivity:}
Let $I = (a, b] \in \mathcal{C}$. Suppose $I = \bigsqcup_{k=1}^n I_k$ is a finite disjoint union of intervals $I_k = (a_k, b_k] \in \mathcal{C}$.
We assume standard ordering $a = a_1 < b_1 = a_2 < \dots < b_n = b$.
\[ \sum_{k=1}^n \lambda(I_k) = \sum_{k=1}^n (b_k - a_k) = (b_1 - a_1) + (b_2 - a_2) + \dots + (b_n - a_n) \]
This is a telescoping sum:
\[ = -a_1 + (b_1 - a_2) + \dots + b_n = b_n - a_1 = b - a = \lambda(I) \]
This argument holds for extended real numbers as well (if the total length is infinite, the sum is infinite).

\textbf{3. Countable Additivity:}
Let $I = \bigsqcup_{n=1}^\infty I_n$ where $I, I_n \in \mathcal{C}$. We need to show $\lambda(I) = \sum_{n=1}^\infty \lambda(I_n)$.
From finite additivity, for any $N$: $\sum_{n=1}^N \lambda(I_n) = \lambda(\bigcup_{n=1}^N I_n) \le \lambda(I)$. Taking $N \to \infty$, we get $\sum \lambda(I_n) \le \lambda(I)$.
We must prove the reverse: $\lambda(I) \le \sum_{n=1}^\infty \lambda(I_n)$.

\textit{Case A: $I$ is bounded ($a, b \in \mathbb{R}$).}
We use the Heine-Borel theorem. Fix $\epsilon > 0$.
Consider the compact interval $K = [a+\epsilon/2, b] \subset (a, b]$.
Expand each $I_n = (a_n, b_n]$ to an open interval $J_n = (a_n, b_n + \epsilon/2^n)$.
Clearly, $\bigcup J_n \supseteq \bigcup I_n = I \supseteq K$. Since $K$ is compact, there is a finite subcover $J_{n_1}, \dots, J_{n_m}$.
The length of $K$ is bounded by the sum of the lengths of the subcover:
\[ (b - (a+\epsilon/2)) \le \sum_{j=1}^m \lambda(J_{n_j}) \le \sum_{n=1}^\infty \lambda(J_n) \]
Substituting $\lambda(J_n) = \lambda(I_n) + \epsilon/2^n$:
\[ b - a - \epsilon/2 \le \sum_{n=1}^\infty (\lambda(I_n) + \frac{\epsilon}{2^n}) = \sum_{n=1}^\infty \lambda(I_n) + \epsilon \]
Since this holds for any $\epsilon > 0$, taking $\epsilon \to 0$ yields $\lambda(I) \le \sum \lambda(I_n)$.

\textit{Case B: $I$ is unbounded.}
If $\lambda(I) = \infty$, the inequality $\lambda(I) \le \sum \lambda(I_n)$ holds trivially if the sum diverges. If $I$ is unbounded, the disjoint union $\bigcup I_n$ must cover an unbounded region, implying $\sum \lambda(I_n)$ must diverge to $\infty$.

Thus, $\lambda$ is a measure on $\mathcal{C}$.
\end{solution}

\hrulefill

\begin{problem}
Consider $\Omega = \mathbb{N}$ and $\mathcal{A} = \mathcal{P}(\mathbb{N})$. Give an example of a set-function on $\mathcal{A}$ which is additive and continuous from above at all $E \in \mathcal{A}$ but not $\sigma$-additive.
\end{problem}

\begin{solution}
Consider the set function $\mu: \mathcal{P}(\mathbb{N}) \to [0, \infty]$ defined by:
\[ \mu(A) = \begin{cases} 0 & \text{if } A \text{ is finite} \\ \infty & \text{if } A \text{ is infinite} \end{cases} \]

\textbf{1. Additivity:}
Let $A, B$ be disjoint sets.
\begin{itemize}
    \item If both are finite, $A \cup B$ is finite. $\mu(A \cup B) = 0 = 0+0$.
    \item If one is infinite (say $A$), $A \cup B$ is infinite. $\mu(A \cup B) = \infty = \infty + \mu(B)$.
    \item If both are infinite, $A \cup B$ is infinite. $\mu(A \cup B) = \infty = \infty + \infty$.
\end{itemize}
Thus $\mu$ is finitely additive.

\textbf{2. Continuity from Above:}
Definition: A set function is continuous from above at $E$ if for any sequence $A_n \downarrow E$ with $\mu(A_k) < \infty$ for some $k$, $\lim_{n \to \infty} \mu(A_n) = \mu(E)$.

Let $A_n \downarrow E$.
\begin{itemize}
    \item If there exists $k$ such that $\mu(A_k) < \infty$, then by definition $A_k$ is a finite set. Since $A_n \subseteq A_k$ for all $n \ge k$, all subsequent sets are finite. Thus $\mu(A_n) = 0$ for all $n \ge k$, and the limit is 0. Since $E \subseteq A_k$, $E$ is finite, so $\mu(E) = 0$. The condition holds.
    \item If $\mu(A_n) = \infty$ for all $n$, the condition ``$\mu(A_k) < \infty$'' is never met, so the continuity property holds vacuously.
\end{itemize}

\textbf{3. Not $\sigma$-additive:}
Let $E_n = \{n\}$ for $n \in \mathbb{N}$. These are disjoint finite sets.
\[ \bigcup_{n=1}^\infty E_n = \mathbb{N} \]
\[ \mu\left(\bigcup_{n=1}^\infty E_n\right) = \mu(\mathbb{N}) = \infty \]
However,
\[ \sum_{n=1}^\infty \mu(E_n) = \sum_{n=1}^\infty 0 = 0 \]
Since $\infty \neq 0$, $\mu$ is not $\sigma$-additive.
\end{solution}

\end{document}